\documentclass[11pt]{article}
\usepackage[a4paper,margin=1in]{geometry}
\usepackage{amsmath,amssymb}
\usepackage{booktabs,longtable}
\usepackage{graphicx}
\usepackage[hyphens]{url}
\usepackage[colorlinks=true,linkcolor=black,urlcolor=blue,citecolor=blue]{hyperref}
\setlength{\parskip}{0.55em}\setlength{\parindent}{0pt}
\providecommand{\tightlist}{\setlength{\itemsep}{0.2em}\setlength{\parskip}{0pt}}

\title{Cover Letter — PLOS ONE}
\author{\textbf{Amin Abaee}\\ Independent Researcher\\[2pt] ORCID: \href{https://orcid.org/0000-0002-0019-1842}{0000-0002-0019-1842}\\ Email: \href{mailto:amin_abaee@ut.ac.ir}{amin_abaee@ut.ac.ir}}
\date{September 18, 2026}
\begin{document}
\maketitle

\section{Cover Letter --- Paper F1}\label{cover-letter-paper-f1}

\textbf{Title:} When a model is not retained: a pre-registered
out-of-sample audit standard for structural modules in empirical
environmental forecasting \textbf{Author:} Amin Abaee (Independent
Researcher) \textbf{Journal:} \emph{PLOS ONE} \textbf{Manuscript
carrier:} \texttt{paperF1\_retention\_framework\_v30\_humanized} (typst
stamp \texttt{se30-c1375ed}), with Supplement S1 v30 \textbf{Keywords:}
forecast evaluation; model retention; pre-registration; negative
results; environmental forecasting; stock assessment; groundwater;
operating characteristics

\begin{center}\rule{0.5\linewidth}{0.5pt}\end{center}

Dear Editor,

I am pleased to submit my manuscript, \textbf{``When a model is not
retained: a pre-registered out-of-sample audit standard for structural
modules in empirical environmental forecasting,''} for consideration in
\emph{PLOS ONE}.

\textbf{What the paper does.} Environmental forecasting
practice---across fisheries, water resources, and beyond---builds
structural modules by default, yet there is no pre-registered standard
that says when added structure has earned its place out-of-sample. This
paper supplies and demonstrates one: a frozen retention rule in which a
candidate causal module is retained only if it reduces rolling-origin
RMSE against both last-value persistence and the next-simpler rung by
more than a stated 5\% tie band, at both one-step and multi-step
horizons, with information availability documented per input at each
forecast origin. The rule is applied to two unrelated,
management-demanding empirical cases---Northern cod biomass (NAFO 2J3KL,
two assessment series) and Edwards Aquifer head at index well J-17
(1934--2023)---with per-origin forecast files archived so every verdict
can be recomputed.

\textbf{Why it matters.} The two worked examples reach empty retained
sets for opposite reasons: on cod, the closest structural approach
misses the band at +9.01\% (five years) and +17.09\% (one year), never
approaching retention, while on Edwards the best-scoring module passes
the predictive gates but is declined on pre-specified class grounds (it
collapses to an affine AR(1) map). The paper then measures the
instrument itself against known ground truth, because a negative verdict
is informative only if the gate has teeth where truth is strong. On
10,000 pre-registered simulation replicates the rule is highly specific
(0.978 under a persistence null), powerful only where signal is strong
(0.955--0.960 at collapse-window parameters), and weak against
stock-flow and depensatory dynamics---power is low for three of four
in-class processes, and the paper shows the cause is parameter
identification (likelihood flatness and compensation), not gate
conservatism: removing all gates adds almost no power. Out-of-class
truths (time-varying productivity, observation error) produce false
retention at 0.633--0.933, delimiting what non-retention can and cannot
certify. A comparison against rule variants on the same archived
replicates (information criterion: 0.509/0.992/0.615; scaled-error gate:
0.651/0.675/0.945; the rule as stated: 0.376/0.978/0.835) locates the
trade the pre-registered rule makes, and the verdicts are shown
band-invariant, with a simulation-calibrated band (power ≥ 0.80,
specificity ≥ 0.90) registered exclusively for future applications---it
never re-opens the verdicts reported here.

\textbf{Why it is appropriate for \emph{PLOS ONE}.} \emph{PLOS ONE}
evaluates rigour and methodology rather than perceived impact, and
explicitly welcomes well-executed negative results---exactly what this
paper is: a methodological standard proven out on two hard real
datasets, calibrated against known synthetic truth, every intermediate
classification reproducible from registered files. The contribution is a
certified instrument for auditing structural complexity claims in any
forecast-driven environmental discipline, together with its measured
limits: which failure modes it detects, which it cannot, and where a
stronger band is licensed only prospectively. That
combination---pre-registration, conflict-free negative certificates
scoped to estimator and series, byte-level reproducibility, and honest
power accounting---sits at the centre of the journal's reproducibility
and methodology remit.

\textbf{Fit and originality.} To my knowledge this is the first
pre-registered forecast-retention audit standard applied jointly across
fisheries and groundwater with archived per-origin replicates,
likelihood-ratio evidential weights per generating mechanism, and a
calibrated prospective band registered before use. The paper's claims
are conservative throughout: it makes no sustainability statement about
either system, draws no unlicensed physical inference from predictive
retention, and states precisely which conclusions are licensed, which
are open, and which are ruled out.

\textbf{Declarations.} All input data, analysis campaigns, frozen result
files, the 10,000-replicate archives, and the independent
band-calibration audits are archived in a public repository
(https://github.com/MIKEAA2020/general-sustainability) and regenerate
deterministically, byte for byte. The work is original, is not under
consideration elsewhere, and I have no competing interests. Funding:
none received.

Thank you for your consideration.

Yours sincerely,

\textbf{Amin Abaee} Independent Researcher ORCID:
\href{https://orcid.org/0000-0002-0019-1842}{0000-0002-0019-1842} Email:
\href{mailto:amin_abaee@ut.ac.ir}{\nolinkurl{amin\_abaee@ut.ac.ir}}

Date: 18 September 2026

\end{document}
